\section{The Grand Canonical Ensemble}

\paragraph{Adding a second constraint: particle exchange.}
Likewise, the \textcolor{blue}{grand canonical ensemble} adds a second
constrained average, the particle number, so the system exchanges both energy
and particles with its surroundings.
\textcolor{red}{Here, $\beta$ and $\gamma$ denote the multipliers of the
constraints $\langle U\rangle=\bar U$ and $\langle N\rangle=\bar N$, and
$U_i$ and $N_i$ denote the energy and particle number of microstate $i$.}
The stationarity conditions of the Lagrangian are
\begin{align*}
    \frac{\partial\mathcal{L}}{\partial p_i}
    &= -\ln p_i-1-\alpha-\beta U_i-\gamma N_i=0,
    &\frac{\partial\mathcal{L}}{\partial\alpha}
    &= -\Big(\sum_i p_i-1\Big)=0,\\
    \frac{\partial\mathcal{L}}{\partial\beta}
    &= -\Big(\sum_i U_i p_i-\bar U\Big)=0,
    &\frac{\partial\mathcal{L}}{\partial\gamma}
    &= -\Big(\sum_i N_i p_i-\bar N\Big)=0.
\end{align*}

\paragraph{The grand canonical distribution.}
The first condition gives $p_i=e^{-1-\alpha-\beta U_i-\gamma N_i}$.
Normalization fixes $\alpha$, and the remaining conditions fix the averages:
\begin{align*}
    e^{1+\alpha}=\sum_i e^{-\beta U_i-\gamma N_i},
    \qquad
    \bar U=\frac{\sum_i U_i\,e^{-\beta U_i-\gamma N_i}}{\sum_i e^{-\beta U_i-\gamma N_i}},
    \qquad
    \bar N=\frac{\sum_i N_i\,e^{-\beta U_i-\gamma N_i}}{\sum_i e^{-\beta U_i-\gamma N_i}}.
\end{align*}
As $\beta$ is the control parameter of the canonical ensemble, $\beta$ and
$\gamma$ are the free parameters here, and $\alpha$ is a normalizing factor
that depends on both. Defining the \textcolor{blue}{grand canonical partition
function}
\begin{equation}
    \boxed{\Xi(\beta,\gamma)\equiv\sum_i e^{-\beta U_i-\gamma N_i},}
\end{equation}
the probability of a microstate with energy $U_i$ and $N_i$ particles is
\begin{equation}
    \boxed{p_i=\frac{e^{-\beta U_i-\gamma N_i}}{\Xi(\beta,\gamma)}.}
\end{equation}

\paragraph{Identifying $\gamma$ with the chemical potential.}
The entropy of this distribution is
\begin{align*}
    \frac{S}{k_B}=-\sum_i p_i\ln p_i
    =\beta\langle U\rangle+\gamma\langle N\rangle+\ln\Xi.
\end{align*}
Differentiating and using $\partial\ln\Xi/\partial\beta=-\langle U\rangle$ and
$\partial\ln\Xi/\partial\gamma=-\langle N\rangle$, the terms from the implicit
dependence cancel, as in the canonical case:
\begin{align*}
    dS=k_B\left(\beta\,d\langle U\rangle+\gamma\,d\langle N\rangle\right),
    \qquad\text{so}\qquad
    \gamma=\frac{1}{k_B}\frac{\partial S}{\partial\langle N\rangle}\bigg|_{\langle U\rangle,V}.
\end{align*}
The first law gives $\partial S/\partial N|_{U,V}=-\mu/T$, hence
$\gamma=-\beta\mu$. Like $\beta=1/(k_BT)$, this is the correspondence with
classical thermodynamics. \textcolor{red}{From here on, we use $(\beta,\mu)$ in
place of $(\beta,\gamma)$.} The results become
\begin{equation}
    \boxed{\Xi(\beta,\mu)=\sum_i e^{-\beta(U_i-\mu N_i)},
    \qquad
    p_i=\frac{e^{-\beta(U_i-\mu N_i)}}{\Xi(\beta,\mu)},}
\end{equation}
and the entropy is
\begin{equation}
    \boxed{S=k_B\left[\beta\left(\langle U\rangle-\mu\langle N\rangle\right)+\ln\Xi\right].}
\end{equation}

\paragraph{Energy and particle number are correlated.}
Exchanging particles also exchanges energy, as the zeroth-law discussion
showed. We therefore expect $\langle U\rangle$ alone to be a poor function of
$\beta$ at fixed $\mu$. To see this, we define the
\textcolor{blue}{covariance} of two observables $X$ and $Y$,
\begin{equation}
    \boxed{\mathrm{Cov}(X,Y)\equiv\langle XY\rangle-\langle X\rangle\langle Y\rangle,
    \qquad\mathrm{Var}(X)=\mathrm{Cov}(X,X).}
\end{equation}
Let $w_i=U_i-\mu N_i$. For any observable $X$, at fixed $V$ and $\mu$,
\begin{align*}
    \frac{\partial\langle X\rangle}{\partial\beta}\bigg|_{V,\mu}
    &=\frac{\partial}{\partial\beta}\left(\frac{\sum_i X_i e^{-\beta w_i}}{\Xi}\right)\\
    &=-\frac{\sum_i X_i w_i e^{-\beta w_i}}{\Xi}
    +\frac{\left(\sum_i X_i e^{-\beta w_i}\right)\left(\sum_j w_j e^{-\beta w_j}\right)}{\Xi^2}\\
    &=-\langle Xw\rangle+\langle X\rangle\langle w\rangle
    =-\mathrm{Cov}(X,w).
\end{align*}
For $X=U$ and $X=w$, this gives
\begin{equation}
    \boxed{\frac{\partial\langle U\rangle}{\partial\beta}\bigg|_{V,\mu}
    =-\mathrm{Var}(U)+\mu\,\mathrm{Cov}(U,N).}
\end{equation}
\begin{equation}
    \boxed{\frac{\partial}{\partial\beta}\left(\langle U\rangle-\mu\langle N\rangle\right)\bigg|_{V,\mu}
    =-\mathrm{Var}(U-\mu N).}
\end{equation}
The monotonicity of $\langle U\rangle(\beta)$ depends on the sign of
$\mu\,\mathrm{Cov}(U,N)$. The combination $\langle U\rangle-\mu\langle N\rangle$
always decreases with $\beta$, like $\langle U\rangle$ in the canonical
ensemble. It follows from $\Xi$ in the same way:
\begin{equation}
    \boxed{\langle U\rangle-\mu\langle N\rangle
    =-\frac{\partial\ln\Xi}{\partial\beta}\bigg|_{V,\mu}.}
\end{equation}

\paragraph{Uniqueness of the multipliers.}
The multipliers $(\beta,\gamma)$ are determined uniquely by the constrained
means $(\langle U\rangle,\langle N\rangle)$ whenever the map between them is
injective. Since $\langle U\rangle=-\partial\ln\Xi/\partial\beta$ and
$\langle N\rangle=-\partial\ln\Xi/\partial\gamma$, the Jacobian of this map is
minus the Hessian of $\ln\Xi$. By the covariance formula above, that Hessian
is the covariance matrix of $(U,N)$:
\begin{equation}
    \boxed{\nabla^2_{(\beta,\gamma)}\ln\Xi
    =\begin{pmatrix}\mathrm{Var}(U)&\mathrm{Cov}(U,N)\\
    \mathrm{Cov}(U,N)&\mathrm{Var}(N)\end{pmatrix}\equiv C.}
\end{equation}
A covariance matrix is positive semidefinite. It is positive definite unless
some combination $aU_i+bN_i$ takes the same value in every accessible
microstate. In the generic case, $\ln\Xi$ is strictly convex, so its gradient
map is injective. Each interior pair $(\bar U,\bar N)$ therefore corresponds to
exactly one pair $(\beta,\gamma)$. This extends the uniqueness argument of the
canonical ensemble to two constraints.

\subsection{Particle Number Fluctuation, Thermodynamic Limit, and Equivalence of Ensembles (again)}

\paragraph{Particle number fluctuation.}
The \textcolor{blue}{particle number fluctuation} is the variance of $N$. At
fixed $\beta$ and $V$, the weights depend on $\mu$ through $e^{\beta\mu N_i}$,
so differentiating the mean gives
\begin{align*}
    \frac{\partial\langle N\rangle}{\partial\mu}\bigg|_{T,V}
    =\beta\left(\langle N^2\rangle-\langle N\rangle^2\right)
    =\beta\,\mathrm{Var}(N).
\end{align*}
Hence
\begin{equation}
    \boxed{\mathrm{Var}(N)=k_BT\,\frac{\partial\langle N\rangle}{\partial\mu}\bigg|_{T,V}.}
\end{equation}

\paragraph{The thermodynamic limit.}
The derivative of an extensive variable with respect to an intensive one is
extensive, so $\mathrm{Var}(N)\sim N$ and
\begin{equation}
    \boxed{\frac{\sqrt{\mathrm{Var}(N)}}{\langle N\rangle}\sim\frac{1}{\sqrt{N}}.}
\end{equation}
As for the energy in the canonical ensemble, particle fluctuations are
negligible because of the sheer size of $N$.
\textcolor{red}{We therefore write $N$ in place of $\langle N\rangle$, and $U$
in place of $\langle U\rangle$, from here on.}
For a very large number of particles, the grand canonical ensemble converges
to the microcanonical one. We can then use the microcanonical ensemble to
solve idealized versions of realistic thermodynamic systems analytically.

\subsection{$\mu VT$ Ensemble - Grand Potential}

\paragraph{The grand potential.}
The natural variables of the grand canonical ensemble are $(\mu,V,T)$, hence
the label $\mu VT$ ensemble. As in the canonical ensemble, $T$ controls the
internal energy, and the chemical potential $\mu$ likewise controls the
particle number. The thermodynamic potential is the
\textcolor{blue}{grand potential}, the partial Legendre transform of
$U(S,V,N)$ in $S$ and $N$:
\begin{equation*}
    \Phi(\mu,V,T)=U-\frac{\partial U}{\partial S}\bigg|_{V,N}S
    -\frac{\partial U}{\partial N}\bigg|_{S,V}N,
\end{equation*}
that is,
\begin{equation}
    \boxed{\Phi(\mu,V,T)=U-TS-\mu N=F-\mu N.}
\end{equation}

\paragraph{Conjugate variables of $\Phi$.}
From $dU=T\,dS-P\,dV+\mu\,dN$,
\begin{align*}
    d\Phi=dU-T\,dS-S\,dT-\mu\,dN-N\,d\mu=-S\,dT-P\,dV-N\,d\mu,
\end{align*}
so
\begin{equation}
    \boxed{\frac{\partial\Phi}{\partial T}\bigg|_{V,\mu}=-S,\qquad
    \frac{\partial\Phi}{\partial V}\bigg|_{T,\mu}=-P,\qquad
    \frac{\partial\Phi}{\partial\mu}\bigg|_{T,V}=-N.}
\end{equation}

\paragraph{Minimization of $\Phi$ at fixed $(\mu,V,T)$.}
At fixed $V$ we have $\delta W=0$, so the first law gives
$\delta Q=dU-\mu\,dN$. The Clausius inequality, $\delta Q\leq T\,dS$, then
gives $dU-\mu\,dN\leq T\,dS$, and at fixed $\mu$ and $T$
\begin{equation}
    \boxed{d\Phi=dU-\mu\,dN-T\,dS\leq0.}
\end{equation}
Any spontaneous process at fixed $(\mu,V,T)$ therefore proceeds toward lower
$\Phi$ and halts when $\Phi$ reaches its minimum.

\paragraph{Grand potential and the partition function.}
Using $S=k_B[\beta(U-\mu N)+\ln\Xi]$ and $\beta=1/(k_BT)$,
\begin{align*}
    \Phi=U-\mu N-TS
    =(U-\mu N)-\frac{1}{k_B\beta}\,k_B\left[\beta(U-\mu N)+\ln\Xi\right]
    =-\frac{\ln\Xi}{\beta},
\end{align*}
so
\begin{equation}
    \boxed{\Phi=-k_BT\ln\Xi.}
\end{equation}

\subsection{$NPT$ Ensemble - Gibbs Free Energy}

\paragraph{The Gibbs free energy.}
The $NPT$ ensemble, conventionally the \textcolor{blue}{isothermal-isobaric
ensemble}, is the one most used in experiments: the particle number $N$, the
pressure $P$, and the temperature $T$ are controlled, and the system exchanges
both energy and volume with its surroundings. Like the $\mu VT$ ensemble, it
has two intensive natural variables, here $P$ and $T$. Its thermodynamic
potential is the \textcolor{blue}{Gibbs free energy}, the partial Legendre
transform of $U(S,V,N)$ in $S$ and $V$:
\begin{equation*}
    G(T,P,N)=U-\frac{\partial U}{\partial S}\bigg|_{V,N}S
    -\frac{\partial U}{\partial V}\bigg|_{S,N}V,
\end{equation*}
that is,
\begin{equation}
    \boxed{G(T,P,N)=U-TS+PV=F+PV=H-TS.}
\end{equation}

\paragraph{Conjugate variables of $G$.}
From $dG=-S\,dT+V\,dP+\mu\,dN$,
\begin{equation}
    \boxed{\frac{\partial G}{\partial T}\bigg|_{P,N}=-S,\qquad
    \frac{\partial G}{\partial P}\bigg|_{T,N}=V,\qquad
    \frac{\partial G}{\partial N}\bigg|_{T,P}=\mu.}
\end{equation}

\paragraph{Minimization of $G$ at fixed $(N,P,T)$.}
At fixed $N$, the first law gives $\delta Q=dU+P\,dV$, and the Clausius
inequality gives $dU+P\,dV\leq T\,dS$. At fixed $P$ and $T$,
\begin{equation}
    \boxed{dG=dU-T\,dS+P\,dV\leq0.}
\end{equation}
Any spontaneous process at fixed $(N,P,T)$ therefore proceeds toward lower
$G$ and halts when $G$ reaches its minimum.

\paragraph{The $NPT$ ensemble from MaxEnt.}
The $NPT$ ensemble arises from the same principle as the others. Constrain
the average volume, $\langle V\rangle=\bar V$, with a multiplier $\beta_V$, and
let the microstates include the volume, with a sum over $V$ replaced by an
integral if $V$ is continuous. The general solution gives
$p_i\propto e^{-\beta U_i-\beta_V V}$. As in the earlier cases,
$\beta_V=\frac{1}{k_B}\,\partial S/\partial\langle V\rangle=\beta P$, because
$\partial S/\partial V|_{U,N}=P/T$. With the partition function
\begin{equation*}
    \Gamma(\beta,P)=\sum_i e^{-\beta(U_i+PV_i)},
\end{equation*}
the entropy is $S=k_B[\beta(\langle U\rangle+P\langle V\rangle)+\ln\Gamma]$.
Solving for the potential gives
\begin{equation}
    \boxed{G=U+PV-TS=-k_BT\ln\Gamma.}
\end{equation}
All five ensembles of this chapter therefore follow from entropy maximization
under different sets of constraints.

\subsection{Euler's Homogeneous Function Theorem}

\paragraph{Extensive variables as functions of other variables.}
We have classified variables as extensive or intensive, but not yet studied
how one extensive variable depends on the others. \textcolor{blue}{Euler's
(homogeneous function) theorem} is a general calculus statement for functions
with a scaling property. A function $f(x_1,\ldots,x_n)$ is
\textcolor{blue}{homogeneous of degree $k$} if
\begin{equation}
    \boxed{f(\lambda x_1,\ldots,\lambda x_n)=\lambda^k f(x_1,\ldots,x_n).}
\end{equation}

\paragraph{Euler's theorem.}
If $f$ is differentiable on $\mathbb{R}^n$ and homogeneous of degree $k$, then
\begin{equation}
    \boxed{\sum_{i=1}^n x_i\frac{\partial f}{\partial x_i}=k\,f(x_1,\ldots,x_n).}
\end{equation}

\paragraph{Proof of the theorem.}
Differentiate the defining relation with respect to $\lambda$. With
$y_i=\lambda x_i$, the chain rule gives
\begin{align*}
    \sum_{i=1}^n\frac{\partial f}{\partial y_i}\bigg|_{y=\lambda x}x_i
    =k\lambda^{k-1}f(x_1,\ldots,x_n).
\end{align*}
Setting $\lambda=1$ yields the theorem.

\paragraph{Application to the internal energy.}
The internal energy is homogeneous of degree one:
\begin{equation*}
    U(\lambda S,\lambda V,\lambda N)=\lambda\,U(S,V,N).
\end{equation*}
Euler's theorem then gives
\begin{equation}
    \boxed{U=\frac{\partial U}{\partial S}\bigg|_{V,N}S
    +\frac{\partial U}{\partial V}\bigg|_{S,N}V
    +\frac{\partial U}{\partial N}\bigg|_{S,V}N=TS-PV+\mu N.}
\end{equation}
This is a relation between values, not a statement about the variables of
$U$. The natural variables of $U$ remain $(S,V,N)$, and $(T,P,\mu)$ are
functions of them, computed from the derivatives of $U$. They cannot be
treated as variables of $U$: scaling the system changes $U$, but leaves
the intensive variables $(T,P,\mu)$ unchanged.

\paragraph{Potentials with fewer extensive variables.}
In the $NPT$ ensemble, $N$ is the only extensive natural variable, so scaling
the system scales $G$ with $N$ alone:
\begin{equation*}
    G(\lambda N,P,T)=\lambda\,G(N,P,T).
\end{equation*}
Euler's theorem then gives $G=\mu N$ for a single-component system, the
formula commonly used in chemistry. The same argument applies to every
potential, and we obtain each one by subtracting its transformed terms from
$U=TS-PV+\mu N$:
\begin{align*}
    F&=U-TS=-PV+\mu N, &H&=U+PV=TS+\mu N,\\
    \Phi&=U-TS-\mu N=-PV, &G&=U-TS+PV=\mu N.
\end{align*}
Hence
\begin{equation}
    \boxed{
    \begin{aligned}
        F(N,V,T)&=-PV+\mu N, &H(N,P,S)&=TS+\mu N,\\
        \Phi(\mu,V,T)&=-PV, &G(N,P,T)&=\mu N.
    \end{aligned}}
\end{equation}
Ensembles with one intensive natural variable, $NVT$ and $NPS$, have a
potential that depends on two extensive variables. Ensembles with two
intensive natural variables, $\mu VT$ and $NPT$, have a potential that
depends on one extensive variable only.

\subsection{Multi-Component System}

\paragraph{Generalization to several species.}
So far, we have treated a single species, but realistic systems contain
several. For species $i$, we introduce the particle number $N_i$ and the
chemical potential $\mu_i$. The pairs $(\mu_i,N_i)$ are conjugate, like the
$(\mu,N)$ we have used, and each species adds one degree of freedom.
\textcolor{red}{From here on, every term $\mu N$ in a Legendre transformation
is replaced by $\sum_i\mu_iN_i$.} Euler's theorem becomes
\begin{equation}
    \boxed{U=TS-PV+\sum_i\mu_iN_i,}
\end{equation}
and all earlier results carry over.

\paragraph{Relations between the potentials.}
The potentials are connected by Legendre transformations:
\begin{center}
\begin{tikzpicture}[
    pot/.style={
        draw, rounded corners=4pt, thick,
        minimum width=3.6cm, minimum height=1.6cm,
        inner sep=4pt, align=center, fill=blue!6
    },
    arr/.style={-{Latex[length=3mm]}, thick, shorten >=2pt, shorten <=2pt},
    lbl/.style={fill=white, inner sep=2pt, font=\small}
]

% Nodes: U and F on the central axis, H/G on the right, Phi mirrored on the left
\node[pot] (U)   at (0,5)    {$U(\{N_i\},V,S)$};
\node[pot] (F)   at (0,0)    {$F(\{N_i\},V,T)$};
\node[pot] (H)   at (7,5)    {$H(\{N_i\},P,S)$};
\node[pot] (G)   at (7,0)    {$G(\{N_i\},P,T)$};
\node[pot] (Phi) at (-7,0)   {$\Phi(\{\mu_i\},V,T)$};

% Right-hand square
\draw[arr] (U) -- node[lbl] {$+PV$} (H);
\draw[arr] (U) -- node[lbl] {$-TS$} (F);
\draw[arr] (H) -- node[lbl] {$-TS$} (G);
\draw[arr] (F) -- node[lbl] {$+PV$} (G);
\draw[arr] (U) -- node[lbl, pos=0.55] {$+PV-TS$} (G);

% Left-hand (mirrored) edges
\draw[arr] (F) -- node[lbl] {$-\sum_i \mu_i N_i$} (Phi);
\draw[arr] (U) -- node[lbl, pos=0.5] {$-TS-\sum_i \mu_i N_i$} (Phi);

\end{tikzpicture}
\end{center}

\paragraph{Summary of the ensembles.}
The table lists, for each ensemble, which variables are fixed and which are
obtained by differentiating the potential:
\begin{center}
\renewcommand{\arraystretch}{2.6}

\begin{tabular}{|c|c|c|c|c|c|}
\hline
Ensemble & $\{N_i\}VS$ & $\{N_i\}VT$ & $\{N_i\}PS$ & $\{\mu_i\}VT$ & $\{N_i\}PT$ \\ \hline
\shortstack{Conventional\\name}
  & \shortstack{micro-\\canonical}
  & canonical
  & \shortstack{isoenthalpic-\\isobaric}
  & \shortstack{grand\\canonical}
  & \shortstack{isothermal-\\isobaric} \\ \hline \hline
\shortstack{Thermodynamic\\potential}
  & $U$ & $F$ & $H$ & $\Phi$ & $G$ \\ \hline
$N_i$
  & fixed & fixed & fixed
  & $-\dfrac{\partial \Phi}{\partial \mu_i}$ & fixed \\ \hline
$V$
  & fixed & fixed
  & $\dfrac{\partial H}{\partial P}$ & fixed
  & $\dfrac{\partial G}{\partial P}$ \\ \hline
$S$
  & fixed & $-\dfrac{\partial F}{\partial T}$ & fixed
  & $-\dfrac{\partial \Phi}{\partial T}$
  & $-\dfrac{\partial G}{\partial T}$ \\ \hline
$\mu_i$
  & $\dfrac{\partial U}{\partial N_i}$
  & $\dfrac{\partial F}{\partial N_i}$
  & $\dfrac{\partial H}{\partial N_i}$
  & fixed
  & $\dfrac{\partial G}{\partial N_i}$ \\ \hline
$P$
  & $-\dfrac{\partial U}{\partial V}$
  & $-\dfrac{\partial F}{\partial V}$ & fixed
  & $-\dfrac{\partial \Phi}{\partial V}$ & fixed \\ \hline
$T$
  & $\dfrac{\partial U}{\partial S}$ & fixed
  & $\dfrac{\partial H}{\partial S}$ & fixed & fixed \\
  \hline
\end{tabular}
\end{center}

\subsection{Mixing Entropy and the Gibbs Paradox}

\paragraph{Setting and the two properties of the entropy.}
Let $X=(U,V,N_1,N_2)$ denote the extensive variables and
$S(X)=k_B\ln\Omega(X)$ the microcanonical entropy. We use
two properties, both understood in the thermodynamic limit.
\begin{itemize}
    \item \textbf{(P1) Superadditivity.} $S(X+Y)\ge S(X)+S(Y)$.
    Place two independent systems with variables $X$ and $Y$ next to each
    other. Each pair of microstates, one from each system, is a microstate of
    the combined system with variables $X+Y$, and distinct pairs give
    distinct microstates. Hence $\Omega(X+Y)\ge\Omega(X)\Omega(Y)$, and taking
    $k_B\ln$ gives (P1).
    \item \textbf{(P2) Homogeneity.} $S(\lambda X)=\lambda S(X)$ for
    $\lambda>0$. This is extensivity. For identical particles it holds only
    if $\Omega$ counts indistinguishable particles, that is, with the factor
    $1/N!$ included.
\end{itemize}

\paragraph{Concavity of the entropy along the composition line.}
For $\lambda\in[0,1]$, (P1) and (P2) give
\begin{align*}
    S\big(\lambda X+(1-\lambda)Y\big)
    \ \ge\ S(\lambda X)+S\big((1-\lambda)Y\big)
    = \lambda S(X)+(1-\lambda)S(Y),
\end{align*}
so $S$ is concave in $X$:
\begin{equation}
    \boxed{S\big(\lambda X+(1-\lambda)Y\big)\ \ge\ \lambda S(X)+(1-\lambda)S(Y).}
\end{equation}
Now fix the total energy $2U$, the volume $2V$, and the total particle number
$2N$. The only remaining variable is the composition, which we label by the
number $t$ of species-2 particles:
\begin{equation*}
    g(t)\equiv S(2U,\,2V,\,2N-t,\,t),\qquad t\in[0,2N].
\end{equation*}
The map $t\mapsto(2U,2V,2N-t,t)$ is affine, and a concave function composed
with an affine map is concave, so $g$ is concave. Three points on this line
matter:
\begin{align*}
    g(0)&=S(2U,2V,2N,0) &&\text{(all species 1)},\\
    g(2N)&=S(2U,2V,0,2N) &&\text{(all species 2)},\\
    g(N)&=S(2U,2V,N,N) &&\text{(half and half)}.
\end{align*}
Since $(N,N)=\tfrac12(2N,0)+\tfrac12(0,2N)$, the half-and-half state is the
midpoint of the two pure states.

\paragraph{Entropy of mixing for distinguishable species.}
Take the two systems before the partition is removed,
$X=(U,V,N,0)$ and $Y=(U,V,0,N)$, and define the entropy of mixing
\begin{equation}
    \boxed{\Delta S_{\mathrm{mix}}\equiv S(2U,2V,N,N)-S(U,V,N,0)-S(U,V,0,N)
    =g(N)-S(U,V,N,0)-S(U,V,0,N).}
\end{equation}
\emph{Step 1: concavity at the midpoint.} Applying the concavity inequality
along the line with $\lambda=\tfrac12$,
\begin{align*}
    g(N)\ \ge\ \tfrac12\,g(0)+\tfrac12\,g(2N).
\end{align*}
The right-hand side is the height of the chord joining the two end points,
at its midpoint, so the graph of $g$ does not lie below its chord.

\emph{Step 2: homogeneity at the end points.} By (P2) with $\lambda=2$,
\begin{align*}
    g(0)=S(2U,2V,2N,0)=2\,S(U,V,N,0),\qquad
    g(2N)=S(2U,2V,0,2N)=2\,S(U,V,0,N).
\end{align*}
Inserting these into the midpoint inequality gives
$g(N)\ge S(U,V,N,0)+S(U,V,0,N)$, that is,
\begin{equation}
    \boxed{\Delta S_{\mathrm{mix}}\ \ge 0.}
\end{equation}
The same result follows directly from (P1), since $X+Y=(2U,2V,N,N)$.

\paragraph{When equality holds.}
Equality in the midpoint inequality means that $g(N)$ lies on the chord. A
concave function that touches its chord at an interior point equals the chord
everywhere: if some $g(t_0)$ were strictly above the chord, concavity would
put $g(N)$ strictly above it too. Therefore
\begin{equation}
    \boxed{\Delta S_{\mathrm{mix}}=0
    \iff g \text{ is affine on } [0,2N],
    \qquad
    \Delta S_{\mathrm{mix}}>0
    \iff g \text{ is not affine.}}
\end{equation}
In particular, $\Delta S_{\mathrm{mix}}>0$ whenever $S$ is strictly concave
along the composition line.

\paragraph{Identical particles and the paradox.}
For identical particles, take $X=Y=(U,V,N)$. By (P2),
\begin{align*}
    S(2U,2V,2N)=2S(U,V,N)=S(X)+S(Y),
\end{align*}
so $\Delta S_{\mathrm{mix}}=0$ exactly. In terms of $g$: if the two
``species'' are physically identical, $S$ depends only on $N_1+N_2$, so $g$
is constant, hence affine, and the equality case applies. The difference
between the two cases is whether the composition is a genuine variable of
$S$. For distinguishable species, $g$ is not affine, and removing the
partition increases the entropy by the strictly positive amount
$\Delta S_{\mathrm{mix}}$. For identical particles, $g$ is constant and
$\Delta S_{\mathrm{mix}}=0$. No continuity is expected between the two: they
differ in what the counting $\Omega$ treats as distinct, which is why (P2)
holds for identical particles only with the $1/N!$ included.

\subsection{Gibbs-Duhem Relation}

\paragraph{The intensive variables are not independent.}
The \textcolor{blue}{Gibbs-Duhem relation} states that the intensive
variables $(T,P,\{\mu_i\})$ are not independent. Write the extensive
variables as $x=(S,V,\{N_i\})$ and their conjugates as
$\xi=(T,-P,\{\mu_i\})$, using $\xi$ as in the Legendre transformation.
Let $f$ be any thermodynamic potential, obtained from $U$ by exchanging the
variables in a set $A$ for their conjugates:
\begin{equation*}
    f=U-\sum_{a\in A}\xi_a x_a,
\end{equation*}
so that $f=U$ for $A=\varnothing$, $f=F$ for $A=\{S\}$, $f=G$ for
$A=\{S,V\}$, and so on. Wherever $f$ is differentiable in its natural
variables,
\begin{equation*}
    df=\sum_{a\notin A}\xi_a\,dx_a-\sum_{a\in A}x_a\,d\xi_a.
\end{equation*}
By the Euler relation, $U=\sum_a\xi_a x_a$, so $f$ is also
\begin{equation*}
    f=\sum_{a\notin A}\xi_a x_a,
    \qquad\text{hence}\qquad
    df=\sum_{a\notin A}\xi_a\,dx_a+\sum_{a\notin A}x_a\,d\xi_a.
\end{equation*}
Comparing the two expressions for $df$ gives
\begin{equation}
    \boxed{\sum_a x_a\,d\xi_a=S\,dT-V\,dP+\sum_i N_i\,d\mu_i=0.}
\end{equation}
The result does not depend on the choice of $A$, so it holds in every
ensemble. Of the $n$ intensive variables, only $n-1$ can be varied
independently. 

\paragraph{Why no ensemble fixes all the intensive variables.}
Exchanging all extensive variables for their conjugates, $A=\{S,V,\{N_i\}\}$,
gives the potential
\begin{equation*}
    f=U-TS+PV-\sum_i\mu_iN_i=0.
\end{equation*}
It vanishes identically by the Euler relation. Such a potential would depend
on no extensive variable at all, so nothing would fix the size of the system.
A system that exchanges energy, volume, and particles freely, with all of
$(T,P,\{\mu_i\})$ held fixed, therefore has no ensemble of its own, and
Gibbs-Duhem is the statement of this fact.